\documentclass[]{article}

\usepackage{amsmath}
\usepackage{multirow}
\usepackage{natbib}
\usepackage{graphicx} 


\begin{document}

Our analysis proceeds in several stages, beginning with the identification of distinct employee trajectories using Latent Class Analysis (LCA), followed by an examination of the organizational and psychological factors that predict membership in these classes.

\subsection{Latent class trajectories of AI-induced job security}
To identify unobserved heterogeneity in employees' perceptions of AI's impact on their job security, we conducted an LCA on the joint distribution of current and expected job security. We evaluated models with two to six latent classes, using the Bayesian Information Criterion (BIC) and normalized entropy to determine the optimal solution. As shown in Figure \ref{fig:image0}, the three-class solution provided the best fit to the data, yielding the lowest BIC value (14,592.21) and a high normalized entropy of 0.803, indicating excellent class separation and classification certainty.

\begin{figure}[h!]
    \centering
    \includegraphics[width=0.5\textwidth]{../input_files/plots/step_2_lca_bic_entropy_1776114673.png}
    \caption{Model selection metrics for the Latent Class Analysis (LCA) of AI-induced job security perceptions. The plot displays the Bayesian Information Criterion (BIC; blue line) and normalized entropy (red line) for models with two to six latent classes. The three-class solution was selected as the optimal model, as it exhibits the lowest BIC value (14,592.21) and a high normalized entropy (0.8032), indicating the best balance of model fit and parsimony with clear class separation.}
    \label{fig:image0}
\end{figure}

The three emergent classes represent distinct psychological trajectories:
\begin{itemize}
    \item \textbf{Class 1: "Resiliently Optimistic" (60.9\% of the sample):} This majority class is characterized by a highly positive and stable outlook. Members of this group overwhelmingly report that AI has a positive impact on their current job security and expect this trend to continue or strengthen over the next three years.
    \item \textbf{Class 2: "Stagnant Neutral" (26.5\% of the sample):} This class exhibits a strong anchoring to the status quo. These employees report that AI has had no impact on their job security and largely expect this to remain the case in the future.
    \item \textbf{Class 3: "Anxiously Declining" (12.6\% of the sample):} This minority class perceives AI as a direct and growing threat. Members report a negative current impact on their job security and anticipate that this negative impact will persist or worsen.
\end{itemize}

The conditional response probabilities for each class, detailed in Table \ref{tab:lca_probs}, provide a granular view of these distinct trajectories. For instance, individuals in the "Resiliently Optimistic" class have a cumulative probability of 84.2\% of reporting a positive current impact, while those in the "Anxiously Declining" class have a 92.4\% probability of reporting a negative current impact.

\begin{table*}[t]
\centering
\caption{LCA Class Sizes and Conditional Response Probabilities for the 3-Class Solution.}
\label{tab:lca_probs}
\resizebox{\textwidth}{!}{
\begin{tabular}{l|cccccc|cccccc}
\hline
\hline
& \multicolumn{6}{c|}{\textbf{Current Impact on Job Security}} & \multicolumn{6}{c}{\textbf{Expected Impact on Job Security (3 years)}} \\
\textbf{Latent Class} & Sig. Neg. & Sl. Neg. & No Impact & Sl. Pos. & Sig. Pos. & Not Sure & Sig. Neg. & Sl. Neg. & No Impact & Sl. Pos. & Sig. Pos. & Not Sure \\
\hline
\textbf{1. Resiliently Optimistic} & 0.000 & 0.030 & 0.124 & 0.419 & 0.423 & 0.004 & 0.000 & 0.046 & 0.120 & 0.377 & 0.449 & 0.008 \\
(60.90\%) & & & & & & & & & & & & \\
\textbf{2. Stagnant Neutral} & 0.000 & 0.051 & 0.861 & 0.052 & 0.000 & 0.037 & 0.000 & 0.152 & 0.628 & 0.173 & 0.000 & 0.047 \\
(26.51\%) & & & & & & & & & & & & \\
\textbf{3. Anxiously Declining} & 0.222 & 0.701 & 0.041 & 0.036 & 0.000 & 0.000 & 0.296 & 0.510 & 0.130 & 0.064 & 0.000 & 0.001 \\
(12.59\%) & & & & & & & & & & & & \\
\hline
\end{tabular}
}
\end{table*}

\subsection{Dimensionality reduction of affective dispositions}
To model employees' underlying emotional attitudes towards AI without multicollinearity, we performed an Exploratory Factor Analysis (EFA) on 16 binary items. The analysis, using a tetrachoric correlation matrix and Varimax rotation, successfully extracted two distinct latent factors. As visualized in the heatmap of factor loadings (Figure \ref{fig:image1}), these factors clearly represent "Positive Affect" (loaded by items such as 'optimistic', 'empowered', 'motivated') and "Negative Affect" (loaded by items such as 'worried', 'threatened', 'anxious'). Both factors demonstrated high internal consistency, with McDonald's Omega values of 0.900 for Positive Affect and 0.901 for Negative Affect. The resulting factor scores were used as continuous covariates in subsequent models.

\begin{figure}[h!]
    \centering
    \includegraphics[width=0.5\textwidth]{../input_files/plots/step_3_qhd_factor_loadings_1776114809.png}
    \caption{Heatmap of factor loadings from an Exploratory Factor Analysis on 16 emotional attitude items (QHD) after a Varimax rotation. The analysis reveals two distinct latent factors, clearly delineating affective dispositions. The first factor, "Positive Affect," is strongly loaded by items such as 'Optimistic', 'Curious', and 'Motivated'. The second factor, "Negative Affect," is primarily characterized by high negative loadings from items like 'Uncertain', 'Skeptical', and 'Anxious', confirming the successful reduction of emotional attitudes into two coherent constructs.}
    \label{fig:image1}
\end{figure}

\subsection{Identifying key organizational predictors with elastic net}
To identify the most salient organizational policies and cultural attributes predicting membership in the "Resiliently Optimistic" class, we employed an Elastic Net logistic regression. This method performs automated feature selection, isolating the most impactful predictors from a larger set of correlated variables.

The results, shown in Figure \ref{fig:image2}, reveal that organizational catalysts for optimism are concentrated around incentive structures and employee agency. The strongest positive predictors were non-monetary incentives, such as peer recognition ($\beta = 0.303$) and learning benefits or certifications ($\beta = 0.236$). Structural enablers that provide employees with agency, such as direct involvement in the AI development process ($\beta = 0.129$) and the freedom to choose AI tools ($\beta = 0.124$), were also highly significant. Conversely, the strongest negative predictors were deterrents rooted in risk and uncertainty, namely distrust in AI accuracy ($\beta = -0.158$) and privacy concerns ($\beta = -0.141$). Notably, the stated importance of regular training yielded a coefficient of zero, suggesting that abstract commitments to training, without tangible incentives or participatory frameworks, are insufficient to foster optimism.

\begin{figure}[h!]
    \centering
    \includegraphics[width=0.5\textwidth]{../input_files/plots/step_4_elastic_net_coefficients_1776115025.png}
    \caption{Elastic Net logistic regression coefficients identifying the organizational policies and employee attitudes that predict membership in the "Resiliently Optimistic" class. The analysis highlights that non-monetary incentives, such as peer recognition (QKC\_2) and learning benefits (QKC\_3), are the most potent positive predictors. Structural enablers, including involving employees in the development process (QKB\_2\_11), also strongly increase the likelihood of optimism. Conversely, the strongest deterrents are distrust in AI accuracy (QF\_6) and privacy concerns (QF\_7), indicating that employee agency and tangible incentives are critical for fostering positive job security perceptions, while risk and uncertainty diminish them.}
    \label{fig:image2}
\end{figure}

\subsection{Multinomial logistic regression of class membership}
Building on the feature selection, we estimated a multinomial logistic regression model to evaluate the adjusted odds of membership in the "Resiliently Optimistic" or "Stagnant Neutral" classes relative to the "Anxiously Declining" reference class. The model included key organizational enablers, deterrents, and affective dispositions.

The odds ratios, visualized in Figure \ref{fig:image3}, confirm and extend the Elastic Net findings. Organizational enablers significantly increased the odds of belonging to the "Resiliently Optimistic" class. For example, providing career progression incentives increased the odds by 15.6\% (OR = 1.156), while direct employee involvement in AI development increased them by 11.1\% (OR = 1.111). The most powerful determinant of class membership was the fear of job loss; employees who cited this as a deterrent had their odds of being in the optimistic class reduced by over 82\% (OR = 0.174) compared to those who did not. Privacy concerns also significantly depressed optimism (OR = 0.855).

\begin{figure}[h!]
    \centering
    \includegraphics[width=0.5\textwidth]{../input_files/plots/step_5_forest_plot_odds_ratios_1776115179.png}
    \caption{Forest plot of odds ratios from the multinomial logistic regression, illustrating the predictors of membership in the "Resiliently Optimistic" class relative to the "Anxiously Declining" reference class. The analysis reveals that deterrents such as worry about job loss (QF\_3) and privacy concerns (QF\_7) significantly decrease the odds of an optimistic outlook. In contrast, organizational enablers, including career progression incentives (QKC\_4), learning benefits (QKC\_3), and employee involvement in AI development (QKB\_2\_11), substantially increase these odds. Furthermore, interaction effects demonstrate that granting employees autonomy over tool selection (QKB\_2\_9\_x\_QF\_7) can buffer the negative impact of privacy anxieties.}
    \label{fig:image3}
\end{figure}

The model also revealed a significant interaction effect: the freedom to choose AI tools moderated the negative impact of privacy concerns. When employees had autonomy over their tools, the negative effect of privacy concerns on their odds of being optimistic was significantly weakened (OR for interaction = 1.281).

\subsection{Simulating the impact of organizational policies}
To translate these regression results into practical insights, we conducted a Monte Carlo policy simulation to estimate the predicted probability of an average employee belonging to the "Resiliently Optimistic" class under different organizational scenarios. We specifically contrasted the effects of providing regular training with direct employee involvement in AI development.

The results, shown in Figure \ref{fig:image4}, demonstrate that policies promoting active agency are far more effective than those focused on passive capability building. In the baseline scenario with neither training nor involvement, the predicted probability of being in the optimistic class is approximately 58\%. Introducing training alone provides only a marginal benefit. However, implementing a policy of direct employee involvement in AI development substantially increases this probability to over 65\%. The combination of both policies yields the highest probability, but the primary driver of this increase is employee involvement. This finding underscores that to foster psychological stability, organizations must grant employees an active, participatory role in the technological transition.

\begin{figure}[h!]
    \centering
    \includegraphics[width=0.5\textwidth]{../input_files/plots/step_6_policy_simulation_probabilities_1776115363.png}
    \caption{Predicted probability of membership in the ``Resiliently Optimistic'' class from a Monte Carlo policy simulation under four organizational scenarios, manipulating the presence of regular training and employee involvement in the AI development process. The simulation demonstrates that involving employees substantially elevates the probability of a positive outlook, both with and without training. In contrast, providing training alone yields only a marginal improvement compared to the baseline, underscoring that active employee participation is a more potent catalyst for psychological stability than passive capability building. Error bars represent uncertainty derived from the simulation.}
    \label{fig:image4}
\end{figure}

\subsection{Robustness of findings}
To ensure the validity of our results, we performed several robustness checks. A Variance Inflation Factor (VIF) analysis of the predictors in the multinomial model confirmed the absence of severe multicollinearity, with all main effects exhibiting VIF values well below the conventional threshold of 5.0. Additionally, we conducted a sensitivity analysis by re-estimating the multinomial model on two modified samples: one excluding respondents who answered "Not sure" to the job security questions, and another where "Not sure" responses were imputed with the modal value. The odds ratios for key predictors, such as employee involvement and fear of job loss, remained remarkably stable across all three model specifications. This stability confirms that our findings are robust and not unduly influenced by the handling of uncertain responses.

\end{document}
                